\documentclass[12pt]{article}
\usepackage[T1]{fontenc}
\usepackage[utf8]{inputenc}
\usepackage{lmodern}
\usepackage{textcomp}
\usepackage{enumitem}
\usepackage{geometry}
\geometry{margin=1in}
\begin{document}
\begin{center}
Answer Key - History - Undergraduate Level Assessment 5 \
\small (Answers are ordered exactly as in the question paper.)
\end{center}

\section*{Multiple Choice}
\begin{enumerate}[label=\arabic*.]
\item \textbf{Answer:} D
\\ \textit{Options:} A) the early migrants were not human.; B) the early migrants did not contribute to the genetic pool of modern Native Americans.; C) the early migrants were a separate species.; D) craniometrics is an unreliable method of tracking populations.; E) skull shapes can change dramatically over generations.
\\ \textit{Note:} The difference in skull shapes between the earliest migrants to the New World and modern Native Americans indicates that craniometrics may not accurately reflect population continuity or lineage, thereby questioning its reliability as a method for tracking human populations over time.
\item \textbf{Answer:} A
\\ \textit{Options:} A) Aztecs.; B) Greeks.; C) Sumerians.; D) Chavin.; E) Chinese.
\\ \textit{Note:} correct because the Maya, while developing their own sophisticated writing and mathematical systems, were influenced by earlier Mesoamerican cultures, including the Aztecs, who contributed to the broader tradition of hieroglyphic writing and numerical systems in the region.
\item \textbf{Answer:} C
\\ \textit{Options:} A) the development of language and writing.; B) slowly increased brain size and the first stone tools.; C) a rapid increase in brain size.; D) the transition from bipedalism to quadrupedalism.; E) a decrease in stature and increase in brain size.
\\ \textit{Note:} After 400,000 years ago, hominid evolution is marked by a significant increase in brain size, which is linked to the development of complex social structures, advanced tool use, and enhanced cognitive abilities that distinguish early Homo species from their predecessors.
\item \textbf{Answer:} B
\\ \textit{Options:} A) coins recovered from ancient shipwrecks.; B) dendrochronology.; C) all of the above.; D) ice core sampling.; E) written records.
\\ \textit{Note:} Dendrochronology provides a precise and annually resolved timeline through tree-ring analysis, which helps to calibrate radiocarbon dating results by correlating carbon-14 dates with known chronological events.
\item \textbf{Answer:} A
\\ \textit{Options:} A) during the final minute of the film; B) 90 minutes into the film; C) about half way through the film; D) during the final 30 seconds of the film; E) in the first hour of the film
\\ \textit{Note:} The correct answer is supported by the concept that the timeline of the universe is vast, with human history representing only a minuscule fraction of that time, thus placing the emergence of early humans in the final minute of the film.
\end{enumerate}

\section*{True / False}
\begin{enumerate}[label=\arabic*.]
\setcounter{enumi}{5}
\item \textbf{Answer:} True
\\ \textit{Options:} A) True; B) False
\\ \textit{Note:} According to the passage: a policy or practice in which members of a majority are discriminated against in favor of a historically disadvantaged group or minority
\item \textbf{Answer:} True
\\ \textit{Options:} A) True; B) False
\\ \textit{Note:} According to the passage: 2010 Caribbean Championships
\item \textbf{Answer:} True
\\ \textit{Options:} A) True; B) False
\\ \textit{Note:} According to the passage: 1201
\item \textbf{Answer:} True
\\ \textit{Options:} A) True; B) False
\\ \textit{Note:} According to the passage: Vermont
\item \textbf{Answer:} False
\\ \textit{Options:} A) True; B) False
\\ \textit{Note:} The correct answer is: 220 to 167 BCE
\end{enumerate}

\section*{Long Form Response}
\begin{enumerate}[label=\arabic*.]
\setcounter{enumi}{10}
\item \textbf{Answer:} Augustus
\\ \textit{Note:} Expected answer: Augustus
\item \textbf{Answer:} to reinforce the badly depleted U.S. Eighth Army
\\ \textit{Note:} Expected answer: to reinforce the badly depleted U.S. Eighth Army
\end{enumerate}

\vfill
\noindent\textit{End of Answer Key}
\end{document}